\section*{Chapter \ref{ch:m1:what-a-trading-firm-does} --- What a Trading Firm Does}

\begin{solution}{exo:m1:what-a-trading-firm-does:1}
Spread $25.02 - 24.98 = \$0.04$; mid $\$25.00$; $0.04/25.00 = \qty{16}{\bp}$.
A market order pays the half-spread, $\$0.02$ a share: $300 \times 0.02 =
\$6.00$.
\end{solution}

\begin{solution}{exo:m1:what-a-trading-firm-does:2}
$2\,000\,000 \times \$0.001 = \$2\,000$ a day, and $\$504\,000$ a year: a
tenth of a cent is not a business at this volume, which is the point of
\cref{prop:m1:what-a-trading-firm-does:breakeven}.
\end{solution}

\begin{solution}{exo:m1:what-a-trading-firm-does:3}
(a) \omterm{def:m1:what-a-trading-firm-does:principal-agent}{Agent}; fee. (b) \omterm{def:m1:what-a-trading-firm-does:principal-agent}{Principal}; spread (the discount), with position risk.
(c) \omterm{def:m1:what-a-trading-firm-does:principal-agent}{Agent} in the economic sense --- the client owns the portfolio; fee.
(d) \omterm{def:m1:what-a-trading-firm-does:principal-agent}{Principal}; spread. (e) \omterm{def:m1:what-a-trading-firm-does:principal-agent}{Principal} for the fund's investors; \omterm{def:m1:what-a-trading-firm-does:alpha}{alpha} --- and,
if the model merely loads on market risk, risk premium.
\end{solution}

\begin{solution}{exo:m1:what-a-trading-firm-does:4}
Gross profit $0.12 \times 500 = \$60$ million. Management fee \$10 million;
performance fee $0.20 \times (60 - 10) = \$10$ million. Investors keep \$40
million: a net return of 8\%. The manager received $20/60$, one third of the
gross profit.
\end{solution}

\begin{solution}{exo:m1:what-a-trading-firm-does:5}
$c = 1.00 - 0.14 \times 5 + 0.20 - 0.05 = 0.45$ cent. $V^\star = 90\,000 /
0.0045 = 20$ million shares a day.
\end{solution}

\begin{solution}{exo:m1:what-a-trading-firm-does:6}
(a) $467.8 \times 10^6 / (1\,545 \times 10^9) = \qty{3.03}{\bp}$. It is an
upper bound because net trading income also includes the firm's trading in
other asset classes, while the denominator counts exchange-traded products
only. (b) $467.8/609 = \text{\euro}0.77$ million per employee. (c)
$50 \times 3.03 \times 10^{-4} = \text{\euro}0.015$: a cent and a half.
\end{solution}

\begin{solution}{exo:m1:what-a-trading-firm-does:7}
(a) Spread earned \$20\,000; position P\&L $-\$9\,173$; net \$14\,427 (after
\$4\,000 of rebates and \$400 of clearing). (b) The dispersion comes from
inventory: the \omterm{def:m1:what-a-trading-firm-does:market-maker}{market maker}'s position is a random walk of $\pm$ one order
size per trade, and that position is exposed to every later move of the mid,
informed jumps included. Quadrupling the order size while dividing the number
of orders by four keeps the volume but makes the inventory walk
$4\sqrt{1/4} = 2$ times larger at any given time of day: the standard
deviation of the daily result roughly doubles, with the same mean.
\end{solution}

\begin{solution}{exo:m1:what-a-trading-firm-does:8}
First, market orders \emph{pay} the half-spread; only resting orders earn it,
and they are not guaranteed to fill: the strategy as described loses 6~basis
points per round trip. Second, even a passive strategy does not keep the
spread: \omterm{def:m1:what-a-trading-firm-does:adverse-selection}{adverse selection} returns a large part of it
(\cref{prop:m1:what-a-trading-firm-does:capture}), and fills are more likely
exactly when they are unprofitable. With ten round trips a day at market the
likely figure is about $-60$ basis points a day before fees.
\end{solution}

\begin{solution}{pb:m1:what-a-trading-firm-does:1}
\textbf{1.} 1 cent; $0.01/40 = \qty{2.5}{\bp}$.
\textbf{2.} $0.12 \times 5 = 0.60$ cent.
\textbf{3.} $c = 1.00 - 0.60 + 0.20 - 0.02 = 0.58$ cent.
\textbf{4.} $6 \times 10^6 \times 0.0058 = \$34\,800$.
\textbf{5.} $\$8\,769\,600$.
\textbf{6.} $4.9 + 2.1 + 0.56 = \$7.56$ million; \$30\,000 a day.
\textbf{7.} \$4\,800 a day; \$1\,209\,600 a year.
\textbf{8.} $30\,000/34\,800 = 86.2\%$.
\textbf{9.} $V^\star = 30\,000/0.0058 = 5\,172\,414$ shares: the firm runs 16\%
above break-even.
\textbf{10.} $c = 1.00 - 0.80 + 0.18 = 0.38$ cent; $V^\star = 7\,894\,737$.
\textbf{11.} $6 \times 10^6 \times 0.0038 - 30\,000 = -\$7\,200$ a day.
\textbf{12.} A half-spread of 1.2 cents, a spread of 2.4 cents, which cannot
be quoted. Quad Street can quote 3 cents and lose most of its fills to
competitors still at 2, or stay at 2 and quote only when its own signals say
the flow is benign: selecting \emph{when} to be in the market is the only
continuous control it has.
\textbf{13.} $c = 0.53$ cent; $V^\star = 5\,660\,377$.
\textbf{14.} $\Phi(-4\,800/22\,000) = \Phi(-0.218) = 0.41$.
\textbf{15.} $0.218 \times \sqrt{252} = 3.46$.
\textbf{16.} $\Phi(-3.46) = 0.027\%$: about one year in 3\,700.
\textbf{17.} $1\,209\,600/5\,000\,000 = 24.2\%$.
\textbf{18.} $(69\,600 - 30\,000)/4\,800 = 8.25$.
\textbf{19.} Volatile markets bring volume and wider spreads, and with fixed
costs almost all of the extra revenue is profit (question 18). They also
bring larger moves against inventory and more informed flow, so the standard
deviation of the daily result and the adverse-selection term rise with them.
\textbf{20.} Quad Street survives above about 5.2 million shares a day, and
the three numbers to watch are its volume, its \omterm{prop:m1:what-a-trading-firm-does:capture}{net capture} per share
(above all the adverse-selection part) and its cost per day.
\end{solution}

\begin{solution}{iq:m1:what-a-trading-firm-does:1}
It earns the half-spread on each fill, plus any exchange rebate, by selling
immediacy to those who want to trade now. It loses in two ways: \omterm{def:m1:what-a-trading-firm-does:adverse-selection}{adverse
selection} --- the fills it gets are disproportionately those just before the
price moves against it --- and inventory risk on the position it accumulates
in the meantime. It pays fixed technology and people costs, so profit is a
small \omterm{prop:m1:what-a-trading-firm-does:capture}{net capture} times a large volume, less those costs.
\iqlookfor{both losses named without prompting, and the sense that the spread
is a gross figure, not the profit.}
\end{solution}

\begin{solution}{iq:m1:what-a-trading-firm-does:2}
A proprietary firm trades its owners' money and earns its trading profit; a
\omterm{def:m1:what-a-trading-firm-does:hedge-fund}{hedge fund} trades investors' money and earns fees on assets and on profits.
Consequences: the fund wants \emph{capacity} (fees scale with assets) and
smooth monthly returns that investors will tolerate; the proprietary firm
wants return on a small capital base and can run strategies with little
capacity, more volatile results, and no obligation to explain them to anyone.
\iqlookfor{incentives derived from who owns the capital, not a list of firm
names.}
\end{solution}

\begin{solution}{iq:m1:what-a-trading-firm-does:3}
You bought 100 at 49.99 and the mid is now 49.97: $-\$2$. Split: $+\$1$ of
spread earned (you bought 1 cent below the then mid of 50.00) and $-\$3$ from
the 3-cent fall of the mid afterwards. The second part, averaged over many
fills, is your \omterm{def:m1:what-a-trading-firm-does:adverse-selection}{adverse selection}.
\iqlookfor{marking at the mid, and the decomposition offered unprompted.}
\end{solution}

\begin{solution}{iq:m1:what-a-trading-firm-does:4}
One basis point of \$2 billion is \$200\,000 a day; times about 250 days,
\$50 million a year. It is plausible: the published figures of listed \omterm{def:m1:what-a-trading-firm-does:market-maker}{market
makers} imply a few basis points of value traded at most, and much less in
the most liquid products.
\iqlookfor{fast arithmetic ($2 \times 10^9 \times 10^{-4}$) and a sanity
check against something known.}
\end{solution}

\begin{solution}{iq:m1:what-a-trading-firm-does:5}
As \omterm{def:m1:what-a-trading-firm-does:principal-agent}{agent}: work the order over several days with an execution algorithm for a
commission; the client keeps the price risk and the information-leakage
risk. As \omterm{def:m1:what-a-trading-firm-does:principal-agent}{principal}: bid for the block at a discount to the last price; the
desk takes the risk and is paid by the discount, which must cover the
expected cost of unwinding ten days of volume plus the volatility over that
period. Hybrids exist: a guaranteed benchmark price for part, agency for the
rest.
\iqlookfor{the risk transfer stated explicitly, and a discount that grows
with size and volatility.}
\end{solution}

\begin{solution}{iq:m1:what-a-trading-firm-does:6}
(i) Spreads widened \emph{because} flow became more informed or more
volatile: \omterm{def:m1:what-a-trading-firm-does:adverse-selection}{adverse selection} per share rose by more than the half-spread.
(ii) You were wider than competitors, so you were filled only when everyone
else had pulled away --- the worst fills --- and your volume was kept up by
exactly those. Measurement: the markout of your fills (the mid a few seconds
and minutes after each fill, signed by your side), split by whether you were
alone at the best price.
\iqlookfor{the idea that fills are selected, and a concrete measurement
rather than a story.}
\end{solution}
